\documentclass[../../main.tex]{subfiles}% 注意这里的文件路径不能用 ./main.tex ,否则用latexmk编译子文件会报错
\graphicspath{{\subfix{./image/}}{\subfix{../../image/}}} % 指定图片引用路径,后续可以直接使用图片文件名
% 如果图片存放在上级目录,则使用 ../../image/ 作为备用路径

% 例如:
% \begin{figure}[H]
% \centering
% \includegraphics[width=0.3\textwidth]{图.png}
% \caption{}
% \label{figure:图}
% \end{figure}
% 注意:上述\label{}一定要放在\caption{}之后,否则引用图片序号会只会显示??.

\begin{document}

\section{草稿一}

\begin{theorem}[Kirszbraun扩张定理, 实值情形]\label{theorem:实值-Kirszbraun扩张定理}
设 $A\subseteq\mathbb{R}^n$ 是任意非空子集, $f:A\to\mathbb{R}$ 是 $L$-Lipschitz 函数, 即对任意 $x,y\in A$,
\begin{align*}
|f(x)-f(y)|\leqslant L|x-y|.
\end{align*}
则存在 $L$-Lipschitz 函数 $F:\mathbb{R}^n\to\mathbb{R}$, 使得 $F|_A=f$, 且对任意 $x,y\in\mathbb{R}^n$,
\begin{align*}
|F(x)-F(y)|\leqslant L|x-y|.
\end{align*}
\end{theorem}
\begin{remark}
一般情形下, Kirszbraun 定理对向量值函数 $f:A\to\mathbb{R}^m$ 也成立, 即存在 $L$-Lipschitz 扩张 $F:\mathbb{R}^n\to\mathbb{R}^m$. 其证明需要更复杂的几何或泛函分析方法.
\end{remark}
\begin{proof}

定义
\begin{align*}
F(x)\triangleq\inf\limits_{a\in A}\left\{f(a)+L|x-a|\right\},\quad x\in\mathbb{R}^n.
\end{align*}
首先证明 $F$ 是 $L$-Lipschitz 的. 任取 $x,y\in\mathbb{R}^n$. 对任意 $a\in A$, 由三角不等式得
\begin{align*}
f(a)+L|x-a|\leqslant f(a)+L|y-a|+L|x-y|.
\end{align*}
对 $a\in A$ 取下确界, 得
\begin{align*}
F(x)\leqslant F(y)+L|x-y|.
\end{align*}
交换 $x$ 与 $y$ 的位置, 同理可得
\begin{align*}
F(y)\leqslant F(x)+L|x-y|.
\end{align*}
因此
\begin{align*}
|F(x)-F(y)|\leqslant L|x-y|.
\end{align*}

其次证明 $F|_A=f$. 任取 $b\in A$. 一方面, 对任意 $a\in A$, 由 $f$ 的 $L$-Lipschitz性得
\begin{align*}
f(a)+L|b-a|\geqslant f(b).
\end{align*}
对 $a\in A$ 取下确界, 得
\begin{align*}
F(b)=\inf\limits_{a\in A}\left\{f(a)+L|b-a|\right\}\geqslant f(b).
\end{align*}
另一方面, 在上式中取 $a=b$, 得
\begin{align*}
F(b)\leqslant f(b)+L|b-b|=f(b).
\end{align*}
故 $F(b)=f(b)$. 由 $b\in A$ 的任意性, $F|_A=f$. 证毕.
\end{proof}
























\end{document}